% Source: https://artofproblemsolving.com/community/c3h1747765_cool_thing
% Stable identity: aops-topic-1747765-q01
The series $\sum\dfrac{1}{k^s}$ converges only for $s>1$, so the formula defines $\zeta(s)$ only there. For $s=-1$ the series is $1+2+3+\cdots$, whose partial sums are
\[
1+2+\cdots+N=\dfrac{N(N+1)}{2},
\]
which grow without bound. So $\sum_{k\ge1}k$ diverges: it is not a number, and it cannot equal $\zeta(-1)$ or anything else.

The first equality in the argument substitutes $s=-1$ into a formula that is not valid at $s=-1$. The value $\zeta(-1)$ is defined in a different way (by extending $\zeta$ from $s>1$ to other $s$, a process called analytic continuation), and with that definition $\zeta(-1)=-\dfrac{1}{12}$; but this does not mean that $1+2+3+\cdots$ converges to $-\dfrac{1}{12}$. $\boxed{\text{The series formula for }\zeta(s)\text{ does not hold at }s=-1\text{; the series diverges.}}$
